\section{Empirical Characterization of Constructive and Destructive Regimes}
\label{sec:gd}

We first isolate curvature with full-batch GD, then introduce sampling while preserving explicit reference objectives. Table~\ref{tab:setup} summarizes the three cases. All diagnostics cover the complete trainable parameter vector. Training uses FP32 and reductions use FP64, with TF32, data augmentation, and dropout disabled. Each full-batch model has one shared initialization across its three rates; the sampling experiment uses three initialization blocks. These are controlled dynamics studies, not matched-compute architecture benchmarks.

\begin{table*}[!t]
\centering\small
\caption{Objectives and coverage. The complete training objective supplies every full-batch GD/Adam update. Representation probes are fixed diagnostic subsets; they do not replace that objective. GD runs have 3,000 updates; the Adam overview has 1,024.}
\label{tab:setup}
\setlength{\tabcolsep}{5pt}
\begin{tabular}{@{}l>{\raggedright\arraybackslash}p{2.0in}>{\raggedright\arraybackslash}p{1.7in}>{\raggedright\arraybackslash}p{1.85in}@{}}
\toprule
 & CIFAR-10 / MLP & CIFAR-10 / ResNet-32 & WikiText-2 / Transformer \\
\midrule
Architecture & Two tanh hidden layers, width 200 & Three residual stages & Two layers, two heads; widths 200 \\
Objective & Half squared error; sum over ten outputs, mean over images & Mean cross-entropy; 50 fixed BatchNorm groups & Token-weighted negative log likelihood \\
Training coverage & First 5,000 training images & All 50,000 training images & 59,680 sequences; 2,088,600 valid targets \\
GD rates & $.004,\ .02,\ .04$ & $.025,\ .05,\ .1$ & $.05,\ .1,\ .2$ \\
Displayed Adam rate & $.001$ & $.003$ & $.003$ \\
Representation probe & All 5,000 training images & First intact 1,000-image group & First 64 sequences; 2,240 positions \\
\bottomrule
\end{tabular}
\end{table*}

\begin{figure*}[!t]
\centering
\includegraphics[width=\textwidth]{\evidence/04-gd-sgd/figures/gd_loss_projected_change_spectral_clean.pdf}
\caption{Training loss and curvature under full-batch GD, with one initialization per model. Top: unsmoothed full-objective loss. Bottom: $Q_s=\eta\widehat\kappa_s$ (solid) and $\eta\lambda_{\max}$ (dashed), both smoothed with the same centered 101-step mean; dotted line: two. Spectra measured every 100 steps are linearly interpolated before smoothing, with equally truncated edge windows. Spectral objectives cover all 5,000 MLP images, 5,000/50,000 ResNet images, and 2,500/59,680 Transformer sequences. Training step $s$ counts completed updates; smoothing does not identify exact onset times.}
\label{fig:gd-loss}
\end{figure*}

\subsection{Directional Curvature in Smooth and Oscillatory Training}
Figure~\ref{fig:gd-loss} distinguishes the realized directional response from the sharpest measured curvature. At the lowest MLP and Transformer rates, loss decreases smoothly and $Q_s$ stays near zero. Higher rates produce oscillations with a continuing downward loss trend and late $Q_s$ near two. ResNet enters an oscillatory regime at every tested rate. A near-zero directional estimate can coexist with a nonzero leading eigenvalue: the gradient need not point along the sharpest direction.

The sparse spectra provide geometric context, not an independent full-objective reference in every model. MLP spectra use the complete objective; ResNet and Transformer use the fixed subsets in the caption. Their largest algebraic eigenvalues have relative eigenpair residual at most $.005$. The identical centered smoother removes a display-filter asymmetry between the curves, but cannot turn a retrospective overlay into evidence that one quantity leads another.

\begin{figure*}[!t]
\centering
\includegraphics[width=\textwidth]{\evidence/04-gd-sgd/figures/gd_reversal_norm_retention.pdf}
\caption{Direction, magnitude, and accumulation for the nine GD conditions in Fig.~\ref{fig:gd-loss}. Consecutive-gradient cosine and gradient norm use complete trailing 50-observation means; normalized displacement $R_s$ is unsmoothed and uses all completed updates. Gradient statistics are placed at their evaluation state, while $R_s$ ends at $\theta_s$. Colors identify the same learning rates as in Fig.~\ref{fig:gd-loss}. Low $R_s$ describes motion relative to path length and is interpreted jointly with loss.}
\label{fig:gd-motion}
\end{figure*}

\subsection{Local Reversal and Cumulative Progress}
Seven conditions develop negative adjacent interaction, while the lowest-rate MLP and Transformer retain constructive behavior (Fig.~\ref{fig:gd-motion}). At the highest rates, $R_{3000}$ is $.0257$, $.0374$, and $.0541$ for MLP, ResNet, and Transformer. Yet those conditions achieve each model's lowest observed endpoint loss: $.1539$, $.4898$, and $6.0697$, respectively. Local reversal therefore coexists with greater absolute displacement and better objective progress over these windows.

The exact interaction audit strengthens this observation. Define $E=\sum_s\|g_s\|^2$, $I_{\rm all}=2\sum_{i<j}g_i^\top g_j$, and $I_1=2\sum_s g_{s-1}^\top g_s$, using the gradients applied in the trajectory. All nine conditions have $I_{\rm all}>0$, including the seven with $I_1<0$. Longer-lag reinforcement more than offsets the negative adjacent sum. Supplementary Table S1 supplies the complete endpoint and interaction values. This result directly rejects adjacent-pair dominance for these trajectories; it is not inferred from a small displacement ratio alone.

\begin{figure*}[!t]
\centering
\includegraphics[width=\textwidth]{\evidence/04-gd-sgd/figures/gd_backward_compact.pdf}
\caption{Backward gradient alignment at completed step $s=2256$. Each point is $\cos(g_s,g_{s-k})$ for one fixed endpoint. Top: lags 1--256; bottom: detail over 1--32. Rates are $.04$, $.1$, and $.2$ for MLP, ResNet, and Transformer. These separately recomputed GD continuations use one initialization; points are neither smoothed nor averaged over endpoints. Lag measures backward separation, not training progress.}
\label{fig:gd-lags}
\end{figure*}

Figure~\ref{fig:gd-lags} resolves part of this temporal organization. Gradients can oppose the immediately preceding direction and align with a more distant one. Alignment can strengthen again after becoming small; its envelope is not generally monotone. These fixed-endpoint profiles expose recurrence beyond lag one, while the complete all-pair audit quantifies accumulation. A profile is not an all-pair average or evidence of repeatedly visiting the same minimum.

\subsection{Sampling at Fixed and Evolving States}
The batch-size comparison consists of 546 continuations: 198 MLP, 150 ResNet, and 198 Transformer. Each lasts 200 updates and starts either at initialization or after 2,000 GD updates. Rates are $.02/.1/.2$. Each stochastic condition has eight sampling repetitions per initialization; the full-GD control has one. Figure~\ref{fig:sgd} averages the 199 adjacent pairs within each continuation, then repetitions, then the three initialization means. Shading is their range, not a confidence interval.

\begin{figure*}[!t]
\centering
\includegraphics[width=\textwidth]{\evidence/04-gd-sgd/figures/sgd_batch_size_compact.pdf}
\caption{Sampled-gradient alignment during 200-update continuations from initialization or a 2,000-update GD state. Lines equally weight three initialization means after repetition averaging; shading spans those means. Hollow rightmost markers denote full GD. ResNet sampling preserves complete fixed 1,000-image normalization groups; Transformer batch sizes count prepared sequences. Batch axes are logarithmic. Equal update budgets entail unequal data exposure.}
\label{fig:sgd}
\end{figure*}

The smallest batches weaken trained-state reversal in all three cases. From initialization, however, small MLP batches can yield negative mean cosine despite positive full-GD alignment. Sampling affects which states are visited as well as the coherence of the measured gradients. In the Transformer, batch eight gives cosine near $-.26$ from both starts, but the mean complete-objective loss falls from initialization and rises from the trained state. Cosine alone cannot rank learning progress.

A separate frozen-state MLP experiment isolates sampling at nine saved states from three initializations. Its exact variance and unnormalized noise-product targets agree with measurements to maximum relative discrepancies of $.50\%$ and $.74\%$. The approximate cosine formula has maximum absolute error $.00745$, but misses 21 of 36 conditional stochastic confidence intervals. Small numerical error and an exact identity are therefore different claims. Supplementary Section S2 reports the sampling protocol, uncertainty, and calibration plots.

\subsection{Independent Curvature Validation and Loss Residuals}
\label{subsec:gd-validation}
The smooth tanh MLP permits an independent derivative check. At 1,188 paired states, we compute the starting-point Rayleigh quotient by HVP and compare its product with $\eta$ to $Q_s$. Each HVP uses the preceding sampled gradient as its direction. For batch 32, the median absolute discrepancy is $.7467$; adding the measured noise term $\varepsilon_s$ reduces it to $.0030$ over 288 pairs. Every stochastic batch improves (Supplementary Table S2). The derivative computation is independent of the two-gradient formula, although observations across updates and initialization blocks are dependent. Full-gradient references make the noise correction an offline validation, not a free online correction.

Curvature agreement does not validate a quadratic loss model on every step. Comparing the surrogate in~\eqref{eq:quadratic-loss} with observed incoming GD changes gives 1,325 sign disagreements among 2,999 paired transitions in the high-rate Transformer. The complete nine-condition audit is retained in Supplementary Table S1. Thus the curvature scalar is informative about gradient response even when the finite-step residual prevents reliable loss-sign inference. This distinction becomes essential under Adam, where the update direction itself changes.
